\documentclass[sigconf,nonacm]{acmart}
\usepackage{tikz}
\usetikzlibrary{positioning}
\usepackage{graphicx}
\usepackage{booktabs}
\settopmatter{printacmref=false,printfolios=false}
\pagestyle{plain}\fancyhead{}\fancyfoot{}
\setcopyright{none}
\acmDOI{}
\setlength{\bibsep}{1pt}
% Layout-only adjustments; no text or figure content changed.
\setlength{\textfloatsep}{10pt plus 2pt minus 2pt}
\setlength{\floatsep}{10pt plus 2pt minus 2pt}
\setlength{\intextsep}{14pt plus 2pt minus 2pt}
\begin{document}
\title{Phishing URL Detection: A Brief Review of the Methods and How They Are Tested}
\author{Prabhnoor Singh}\affiliation{\institution{University of Victoria}\city{Victoria}\country{Canada}}\email{prabhnoorsingh@uvic.ca}
\author{Divraj Singh}\affiliation{\institution{University of Victoria}\city{Victoria}\country{Canada}}\email{divrajsingh@uvic.ca}
\author{Dilraj Gulati}\affiliation{\institution{University of Victoria}\city{Victoria}\country{Canada}}\email{dilrajsinghgulati@uvic.ca}
\begin{abstract}
Phishing links trick people into handing over passwords and money. This review traces the problem's history, groups detectors (from blacklists to machine learning on the link itself) by the information they use, and compares representative methods. We also examine how detectors were tested: many are tested on data from the same collection they were trained on, and few check unseen domains. Our project focuses on this evaluation gap. Code and project records: \url{https://github.com/prabhnoob/domain-generalizable-phishing-url-detection}.
\end{abstract}
\keywords{Phishing detection, URL classification, evaluation, literature review}
\maketitle
\section{Introduction}
A phishing detector decides whether a link is safe or a trap. Researchers have used four main approaches: checking lists of known bad links, inspecting webpages, learning patterns from URL text, and looking at connections between URLs. Lists depend on reports of new attacks; page-based methods need to load a page; URL-based methods are fast but may learn surface patterns; and network-based methods need information about related links. Our project studies whether URL detectors still work when tested on websites they have not seen before.
\section{History}
Tricking people by pretending to be someone they trust is far older than computers. The word ``phishing'' draws on the ``ph'' of phone phreaking and was linked to online credential theft in 1996, when attackers stole America Online accounts \cite{gupta}. An early large-scale defence was the blacklist, a list of reported phishing sites checked by browsers. Most blacklists caught fewer than 20\% of new phish at ``hour zero,'' while 63\% of campaigns lasted under two hours \cite{sheng}. Research therefore turned to unreported links. CANTINA used webpage text in 2007 \cite{cantina}; by 2009--2010, machine learning was used for URLs and pages \cite{ma,whittaker}. From 2017, deep learning reduced the need for hand-made features \cite{bahnsen,urlnet}, followed by transformers \cite{urltran} and large language models \cite{llm}. In 2025, over 3.8 million attacks were recorded \cite{apwg}.
\section{Hierarchy of Approaches}
Following the broad categories in Gupta et al. \cite{gupta}, we group methods by the information they use (Figure~\ref{fig:taxonomy}). List-based methods check known URLs. Page-based methods examine text, HTML, or screenshots. URL-based methods read the link, sometimes with host facts. These divide into hand-crafted features and features learned by a model. Network-based methods use connections between URLs. How a method is \emph{tested} is a separate question, not another detection category.
\begin{figure}[t]
\centering\resizebox{\columnwidth}{!}{%
\begin{tikzpicture}[font=\scriptsize,level distance=12mm,
 every node/.style={align=left}, edge from parent/.style={draw,-},
 level 1/.style={sibling distance=33mm},level 2/.style={sibling distance=18mm}]
\node {Phishing detection}
 child {node {List-based\\\cite{sheng}}}
 child {node {Page-based} child {node {Content\\\cite{cantina}}} child {node {Visual\\\cite{visual}}}}
 child {node {URL-based\\(our focus)} child {node {Hand-crafted\\\cite{ma,sahingoz}}} child {node {Learned\\\cite{bahnsen,urlnet,urltran,llm}}}}
 child {node {Network-based\\\cite{kim}}};
\end{tikzpicture}}
\caption{Detection approaches grouped by information used. Evaluation methods are discussed separately.}\label{fig:taxonomy}
\end{figure}
\begin{table*}[t]
\centering\caption{Representative approaches and evaluation methods. U = URL, P = page, N = host/network facts. Results come from different studies and are not directly comparable.}\label{tab:comparison}
\scriptsize\renewcommand{\arraystretch}{1.3}\setlength{\tabcolsep}{3.5pt}
\begin{tabular}{@{}lllllll@{}}\toprule
Category & Example & Uses & Main strength & Main limitation & Test split & Unseen domains\\\midrule
List-based & Blacklists \cite{sheng} & U & Few false alarms & Misses new attacks & -- & --\\
Page: content & CANTINA \cite{cantina} & P,U,N & Uses webpage text & Needs page and search & -- & --\\
Page: visual & VisualPhishNet \cite{visual} & P & Detects copied designs & Needs page rendering & -- & --\\
URL: features & Ma et al. \cite{ma} & U,N & Simple, explainable & Hand-made features & Random URLs & No\\
URL: CNN & URLNet \cite{urlnet} & U & Learns URL patterns & May miss new patterns & Time period & No\\
URL: transformer & URLTran \cite{urltran} & U & Learns URL text & Larger model & Week & No\\
Network-based & Kim et al. \cite{kim} & U,N & Uses linked URLs & Needs network data & -- & --\\
Evaluation & Rashid et al. \cite{rashid} & U & New data source & Domain overlap possible & Dataset & Partial\\
Evaluation & Ahamed et al. \cite{ahamed} & U & Stress tests & Domain overlap remains & Host + dataset & Partial\\
Our project & Two URL baselines & U & Domain + source tests & No new detector & Domain + source & Yes\\\bottomrule
\end{tabular}
\end{table*}
\section{Comparative Analysis}
Table~\ref{tab:comparison} compares the main methods. Results come from different tests, so their scores should not be treated as a direct ranking.
\textbf{Lists and pages.} Blacklists are precise for reported sites but often miss new attacks \cite{sheng}. Whittaker et al. \cite{whittaker} used URLs and webpage content to add sites to Google's blacklist. CANTINA \cite{cantina} uses important words from a page and a web search, with extra checks such as domain age. VisualPhishNet \cite{visual} compares screenshots with known brands. Page-based methods can see more than the URL, but must load the page.
\textbf{URL-based methods.} Ma et al. \cite{ma} used URL tokens and host facts with a classifier; Sahingoz et al. \cite{sahingoz} built language-style URL features for a random forest. These methods depend on features chosen by people. Bahnsen et al. \cite{bahnsen} instead used an LSTM to learn from URL characters. URLNet \cite{urlnet} learned from characters and words using CNNs, while URLTran \cite{urltran} used a transformer. LLMs have also been tested as one-shot classifiers \cite{llm}. Learned models need less manual feature design, but can be larger and harder to run. High scores on familiar data do not guarantee that a model works on new websites \cite{ahamed}.
\textbf{Networks and testing.} Kim et al. \cite{kim} use related URLs to make detection harder to evade, but need connection data. Arp et al. \cite{arp} warn that data leakage can make security models look better than they are. Rashid et al. \cite{rashid} found weaker results when detectors were tested on another dataset. Ahamed et al. \cite{ahamed} tested several models under realistic changes, but 1,387 registrable domains still appeared in both training and testing. This motivates our stricter split.
\begin{figure}[t]
\centering
\begin{tikzpicture}[scale=1.02, every node/.style={font=\scriptsize,align=center}]
% Overlapping circles follow the compact, muted style of the manual.
\fill[green!38,opacity=.72] (-.85,.28) circle (1.35);
\fill[gray!28,opacity=.72] (.85,.28) circle (1.35);
\fill[gray!45,opacity=.65] (0,-.55) circle (1.35);
\draw[gray!65,thin] (-.85,.28) circle (1.35);
\draw[gray!65,thin] (.85,.28) circle (1.35);
\draw[gray!65,thin] (0,-.55) circle (1.35);
% Labels name the information sources; examples sit inside each region.
\node[font=\small\bfseries] at (-1.55,1.80) {URL};
\node[font=\small\bfseries] at (1.55,1.80) {Page};
\node[font=\small\bfseries] at (0,-2.06) {Host / network};
\node at (-1.45,.78) {URLNet\\URLTran\\Our project};
\node at (1.3,.82) {VisualPhishNet};
\node at (-.8,-.5) {Ma et al.\\Kim et al.};
\node at (0,-.05) {CANTINA};
\node at (0,1.12) {Whittaker};
\end{tikzpicture}
\caption{Information used by representative methods. The overlaps show methods that combine sources; our project uses only the URL.}\label{fig:venn}
\end{figure}
\section{Positioning Our Project}
Our contribution is to test whether URL-based phishing detectors work on unseen websites and on data from another source. We focus on URLs because they can be checked without opening the page or collecting a network (Figure~\ref{fig:venn}). We will compare a model using hand-made URL features, inspired by Sahingoz et al. \cite{sahingoz}, with one learning from URL characters, inspired by Bahnsen et al. and URLNet \cite{bahnsen,urlnet}. Both will be trained on PhiUSIIL \cite{phius}. We will compare common evaluation splits with a stricter split that keeps each registrable domain (for example, paypal.com) on only one side. We will also test on another data source. We do not assume either model will perform better; we will report the results of each test.
\noindent\textbf{Team contributions.} Prabhnoor: testing review and Table 1. Divraj: URL and network review and Figure 2. Dilraj: history, Figure 1, list/page review, repository. All: remaining work.
\bibliographystyle{ACM-Reference-Format}
\bibliography{references}
\end{document}
